\section{Ameaças à Validade}
\label{sec:ameacas}

A principal limitação é o tamanho da amostra manual. Foram 60 \textit{templates} de 20 participantes, contra 90 por cenário de IA. Esse tamanho foi limitado pela dificuldade de recrutar voluntários que concluíssem a atividade dentro do período de coleta, não por qualquer exigência de experiência prévia, que não houve. Terraform é uma tecnologia pouco presente nas disciplinas de graduação e no dia a dia da maioria dos estudantes de tecnologia, o que reduziu o número de pessoas dispostas e disponíveis para terminar a atividade manual no prazo. Por causa disso, os resultados do cenário manual são mais sensíveis a casos isolados do que os dos cenários automatizados, que têm 30 execuções cada.

Há também limitações no ambiente de teste. O LocalStack usa uma lista de variáveis de entrada e um roteiro de verificação por \textit{output}, os dois montados a partir dos nomes mais comuns encontrados nos \textit{templates}. Um \textit{template} que usa um nome de variável ou de \textit{output} fora dessa lista, ou que expõe o \textit{output} em um formato diferente, como um mapa com vários identificadores em vez de um único valor, pode não ser testado corretamente nessa etapa, mesmo que o recurso em si esteja certo. O emulador também não implementa os serviços S3 Control e KMS da AWS. Então, \textit{templates} cujo código usa essas APIs, direta ou indiretamente, reprovam a execução funcional por uma limitação do ambiente, não por um erro do \textit{template}. Por fim, a verificação \texttt{CKV2\_AWS\_5} do Checkov sempre reprova nos três cenários, porque ela exige que o Security Group esteja associado a outro recurso, e isso é esperado no desenho deste experimento, que avalia cada Security Group isoladamente.

A padronização foi medida apenas pela presença dos arquivos esperados, sem verificar a aderência às convenções de nomenclatura e \textit{tags} do contexto organizacional. Três participantes manuais enviaram o próprio arquivo de travamento de dependências, e seus \textit{templates} foram executados no LocalStack com o \textit{provider} AWS 6.68.0, e não com o 5.100.0 usado nos demais.

Os resultados refletem um único modelo (Claude, variante \textit{Sonnet}, identificado nos metadados só por esse nome, sem a versão exata) e um único conjunto de \textit{prompts} e documento de contexto organizacional. Não é objetivo deste trabalho comparar modelos de IA entre si, mas sim comparar métodos de criação de \textit{templates}. Por isso, testar outros modelos, versões ou formulações de contexto fica como trabalho futuro.
